\subsection{中继部署与任务时序的联合求解}
\subsubsection{候选悬停位置生成与可服务区间计算}
本节区分题目要求的联合决策空间与实际搜索的有限子域。求解从已保存的23项运输任务和4个中继位置起步，对位置、高度、资源顺序和通信切换进行局部改进；另以小规模运输任务重组检验是否能进一步改善。候选的初始箱组、路线及部分悬停位置继承自补充资料，具体来源与本项目新增计算分列于附录\ref{app:source_provenance}，不把所有初始结构描述为从零构造。

\begin{table}[htbp]\centering\small
\caption{联合求解各阶段的开放变量与固定范围}\label{tab:q3_decision_scope}
\begin{tabularx}{\linewidth}{p{2.4cm}XX}\toprule
阶段 & 允许改变的对象 & 固定内容与输出\\\midrule
位置候选搜索 & 单个中继的水平位置和离地高度 & 暂固定运输模式与其他中继；输出少量候选，尚不宣称完整可行\\
运输重组对照 & 小货箱并集中的组批、机型、访问顺序及选中任务数 & 只在显式给定的有限任务池内选择；未截断的枚举与截断池分开记录\\
联合整数排程 & 候选任务选择、开始时刻、提供方、飞机/电池及中继实体顺序 & 单个任务的路线、单个中继槽的位置已给定；输出可行离散结构\\
连续时间精修 & 运输开始、中继服务开始/结束及联合工期 & 固定路线、位置、提供方、真实资源顺序与充电分段；输出时刻\\
独立验收 & 不再改变计划 & 从完整几何重建区间、端点、能量与资源占用，接受或拒绝候选\\\bottomrule
\end{tabularx}
\end{table}

位置生成在当前点附近按水平网格偏移，并枚举25、50、\ldots、300\,m共12个离地高度。默认搜索半径1200\,m、网格300\,m、每个角色保留10个候选；保存的局部试验还使用75\,m和50\,m尺度，以及面向关键角色的扩大探测，均为有限试验域。先检查DEM覆盖、AGL、回传链路和中继转场，再以任务阶段严格内部、约25\,s间隔的采样估计接入潜力。评分按角色优先考虑出程、回程或二者之和，并以链路裕度破同分；保留水平位置互异的短名单。这里的抽样只可能筛选候选，不能作为整个连续空间不可行的证据。

每个保留点都重新建立完整几何。运输机在一个飞行阶段内位置随时间线性变化；固定端点至该阶段的射线扇面与DEM地形棱柱求交，得到遮挡事件区间，再加入距离门限的解析交点、飞行阶段和服务窗口边界。相邻端点之间的开区间内状态固定，检查内部状态并单独检查全部临界端点，保留孤立擦边事件。这一步才产生式\eqref{eq:q3_cover}使用的可服务关系，不用采样连线替代。

候选搜索先比较完整求解后的工期，再比较同工期能耗；当前实现的接受阈值为工期改善超过0.005\,s，或工期差不超过0.005\,s时能耗改善超过0.005\,kWh。该阈值仅用于有限搜索的数值筛选，不代表实际执行的抗扰动余量。一次角色扫描或预定网格完成后停止，不能宣称无导数全空间搜索已经收敛。

\subsubsection{通信提供方与资源顺序的联合选择}
对有限运输任务池$\mathcal I$取选择变量$x_i$，逐箱覆盖沿用式\eqref{eq:q2_cover}；只重排已有23项任务时固定$x_i=1$。每个中继槽$q$对应一个候选位置，取启用变量$y_q$及式\eqref{eq:q3_cover}的提供方变量$z_{hq}$。需求属于任务$i$时，将覆盖等式右端写为$x_i$，未选择任务不再产生通信需求；$z_{hq}\le y_q$，同时禁止启用不服务任何需求的中继槽。

有限模型通过\emph{任务可接续路径}控制资源数量。对同型飞机，二元变量$v^U_{ij}$表示任务$j$紧接$i$，$o_i^U$表示任务$i$占用一条资源链的起点，满足
\begin{equation}
\begin{aligned}
 o_i^U+\sum_jv^U_{ji}&=x_i,&\sum_jv^U_{ij}&\le x_i,\\
 \sum_{i:m_i=m}o_i^U&\le n_m^U,&
 v^U_{ij}=1&\Rightarrow s_j\ge s_i+p_i .
\end{aligned}
\label{eq:q3_transport_paths}
\end{equation}
电池路径将$p_i$替换为$p_i+c_i$，起点数上限替换为该型号电池库存。正时长的先后约束排除环，故每条路径可以赋予一个实体资源编号。中继实体同样构造路径，接续条件为后项任务开始不早于前项返航加300\,s周转。条件式以二元变量和由有限时间上界计算的大常数实现，不用一个无根据的无限大数。

选定位置后，中继服务开始$a_q$、结束$b_q$以及任务开始$\sigma_q$满足式\eqref{eq:q3_lifecycle}的启用形式：$\sigma_q=a_q-\ell_qy_q$，$\ell_q=180+t^{out}_{p_q}+30$；$b_q\ge a_q$，能耗为转场项、建链项及服务时长项之和。服务覆盖用$z_{hq}$连接运输和中继时间轴。运输硬截止在该阶段另保留30\,s的算法性余量，期望送达仍要求零迟到；30\,s不是新添题设，而是本次搜索采用的保守限制。

当前中继架次试验最多启用6个槽，初始满电能源组件为6组，因此内层整数模型可先让每个启用中继任务使用不同组件，保证存在不复用的组件指派，不需要在此阶段对组件再建组合接续变量。选出任务后，按真实含充电区间重新构造组件复用链，正式方案只需3组；后续连续精修显式保留这3组组件的恢复关系。这一简化依赖“启用任务数不超过初始组件数”，不能直接推广到更多中继架次。

先用HiGHS求有限模型的最小联合工期，再在已得工期增加不超过$10^{-5}$\,s的条件下，允许整数关系继续改变并最小化总能耗。接着固定实际资源链进入下一节的连续精修。当前保存的设置为单线程、整数相对间隙$10^{-7}$、整数可行性容差$10^{-7}$；公开重复优化入口默认90\,s预算，历史各局部试验的实际限时分别保存在状态文件中，不把入口默认值当作所有试验的实际耗时。

\textbf{运输结构和架次数的作用。}除固定23项任务外，程序也对小货箱并集重新生成模式：并集不超过11箱时枚举非空子集、可用机型及路线；模式超过250项时按单位箱数时间、单位箱数能耗和最迟启动余量交错截断。保存对照覆盖S006/S007、S004/S005、S003/S015等组合，联合任务池最大361项，未得到比对应参考更好的已验证运输结构。因此本章主方案与同模型改进前方案的23项箱组、机型和有序路线逐项相同，约87.998\,s的成功改进应归于中继位置与联合时序，而不是运输重组。

\begin{table}[htbp]\centering\small
\caption{不同运输及中继架次试验的实际范围}\label{tab:q3_sortie_scope}
\begin{tabularx}{\linewidth}{p{3.0cm}XX}\toprule
对照 & 搜索范围与结果 & 可作出的结论\\\midrule
运输模式重组 & 小并集与有限路线池，最高361个候选；部分试验限时，未选出更好的正式结构 & 运输架次数可在池内改变；主方案23架次是最终选中结果，不是全局最少证明\\
3个中继架次 & 给定所测4个位置和完成时间上限，有限模型不可行 & 只排除该布局及时间阈值下的三架次方案\\
最多5或6架次 & 8个同位置重复出动槽；普通模型限时未改善 & 未改善或超时不证明所有增加架次的方案无效\\
最多5架次加对称消除 & 同位置槽安全重标号后，有限模型最优仍选4项 & 该有限模型无需第5项；不构成连续位置空间的四架次下界\\\bottomrule
\end{tabularx}
\end{table}

阶段性保存工期依次出现5924.967、5895.017、5866.141、5857.090和5836.970\,s。位置、资源顺序和能耗再优化在不同阶段同时调整，所以该轨迹是实际搜索进展，不是逐因素可加的消融分解。全部保存状态包括110份整数模型结果（91份有限域最优、15份有限域不可行、4份限时）及125份连续精修记录；记录数不能等同于不同有效方案数。

\subsubsection{给定任务结构后的时间安排优化}
\label{sec:q3_fixed_lp}
为使时刻精修可直接复核，本节给出实际使用的固定结构线性模型。输入固定运输箱组、机型、路线，中继位置，通信区间的提供方，以及飞机、电池、中继实体和组件的前后顺序。变量为运输开始$s_i$、中继服务开始$a_q$、服务结束$b_q$和联合最后返航$M$。记$\ell_q=180+t_{p_q}^{out}+30$，则任务开始$\sigma_q=a_q-\ell_q$、返航$e_q=b_q+t_{p_q}^{back}$。

对运输任务，根据箱级要求计算
\begin{equation}
 \bar s_i=\min\left\{\min_{b\in\mathcal B_i}(d_b-\delta_{ib}),\;
 \min_{b\in\mathcal B_i\cap\mathcal B_H}(h_b-\delta_{ib}-30)\right\},
 \quad 0\le s_i\le\bar s_i .
 \label{eq:q3_lp_start_bounds}
\end{equation}
$\mathcal B_H$是存在硬截止的货箱集合，空集最小值取正无穷。已固定给中继$q$的相对需求区间$[\ell_h,u_h]$满足
\begin{equation}
 a_q\le s_i+\ell_h-g_s,\qquad b_q\ge s_i+u_h+g_s,
 \quad a_q\ge\ell_q,\quad b_q\ge a_q .
 \label{eq:q3_lp_coverage}
\end{equation}
实际$g_s=10^{-4}$\,s是数值保护间隔。提供方在进入此模型前按完整可服务关系确定，线性规划本身不再自由改派或移动中继。

固定位置的转场能耗为$E_q^{tr}$，令$w=1.10/3600$\,kWh/s，则
\begin{equation}
 E_q=E_q^{tr}+w(30+b_q-a_q),\qquad 0\le E_q\le0.8B_R,
 \quad B_R=3.2\ \mathrm{kWh}.
 \label{eq:q3_lp_energy}
\end{equation}
充电按名义任务所在的分段固定，并同时施加保持该分段的能量条件，不能只线性化而遗漏分段边界。若$E_q\le0.1B_R$，则$c_q=(3.5\tau_R^{full}/B_R)E_q$；若$E_q\ge0.1B_R$，则
\begin{equation}
 c_q=\frac{0.65\tau_R^{full}}{0.9B_R}E_q
       +\tau_R^{full}\left(0.35-\frac{0.65\times0.1}{0.9}\right).
 \label{eq:q3_lp_charge}
\end{equation}
两段在$E_q=0.1B_R$处连续。该表达由式\eqref{eq:common_recharge}代入$SOC=1-E_q/B_R$得到。

对固定资源链上的相邻任务，分别要求
\begin{equation}
\begin{aligned}
 s_j&\ge s_i+p_i+g_r &&\text{（同一运输机）},\\
 s_j&\ge s_i+p_i+c_i+g_r &&\text{（同一运输电池）},\\
 a_v-\ell_v&\ge b_q+t_{p_q}^{back}+300+g_r &&\text{（同一中继机）},\\
 a_v-\ell_v&\ge b_q+t_{p_q}^{back}+c_q+g_r &&\text{（同一中继组件）}.
\end{aligned}
\label{eq:q3_lp_resource}
\end{equation}
其中$g_r=10^{-4}$\,s，运输$c_i$因任务属性固定而为常数；中继$c_q$由式\eqref{eq:q3_lp_energy}--\eqref{eq:q3_lp_charge}给出仿射函数。目标为
\begin{equation}
 \min M,\qquad s_i+p_i\le M,\quad b_q+t_{p_q}^{back}\le M .
 \label{eq:q3_lp_objective}
\end{equation}
还保留实现中的20000\,s变量上界；当前解远未触及该数值边界。求出$M_{LP}^*$后，再加$M\le M_{LP}^*+10^{-6}$\,s并最小化$\sum_q E_q$。两轮均使用SciPy HiGHS，原始及对偶可行性容差为$10^{-8}$。

该线性规划的工期下界为5836.969928328251\,s，二次节能后的保存值比其大$10^{-6}$\,s。因此可称给定路线、位置、提供方、资源顺序和充电分段的连续时序已闭合；不能称原问题全局最优，也不能把数值保护间隔或求解精度当作实飞时序容错。最终计划仍返回未修改的完整物理与连续通信验证器验收。
